% © 2026 Ing. David Jaroš — CC BY-NC-ND 4.0
%
% This work is licensed under a Creative Commons Attribution-NonCommercial-NoDerivatives
% 4.0 International License (CC BY-NC-ND 4.0).
%
% License History: Earlier drafts (up to v0.3) were released under CC BY 4.0.
% From v0.4 onward, all material is released under CC BY-NC-ND 4.0 to protect
% the integrity of the theoretical work during ongoing academic development.
%
% See LICENSE.md for full license text.

\ifdefined\INCLUDEMODE
\else
  \documentclass[12pt,a4paper]{article}
  \usepackage[a4paper, margin=2.5cm]{geometry}
  \usepackage{amsmath, amssymb, amsthm}
  \usepackage{mathtools}
  \usepackage{hyperref}
  \usepackage{xcolor}
  \usepackage{mdframed}
  \usepackage{booktabs}

  \newtheorem{theorem}{Theorem}[section]
  \newtheorem{lemma}[theorem]{Lemma}
  \newtheorem{proposition}[theorem]{Proposition}
  \newtheorem{corollary}[theorem]{Corollary}
  \theoremstyle{definition}
  \newtheorem{definition}[theorem]{Definition}
  \theoremstyle{remark}
  \newtheorem{remark}[theorem]{Remark}

  \newmdenv[backgroundcolor=yellow!20, linecolor=orange, linewidth=1.5pt]{gapbox}
  \newmdenv[backgroundcolor=green!15, linecolor=green!60!black, linewidth=1.5pt]{resultbox}
  \newmdenv[backgroundcolor=blue!8, linecolor=blue!50, linewidth=1.5pt]{insightbox}

  \title{\textbf{Discrete Symmetry Formalization, Step 1:\\
                 Definitions of C, P, T, and CPT\\
                 in Biquaternion Field Theory}}
  \author{David Jaro\v{s}}
  \date{2026-04-25}

  \begin{document}
  \maketitle

  \begin{abstract}
  We define the discrete symmetry transformations---charge conjugation $C$,
  spatial parity $P$, time reversal $T$, and their composition $CPT$---within
  the biquaternion algebra $\mathcal{B} = \mathbb{C} \otimes \mathbb{H}$ and
  for the fundamental UBT field $\Theta(q,\tau)$ over complex time
  $\tau = t + i\psi$.
  Each transformation is expressed as a composition of the three canonical
  involutions $P_1$ (complex conjugation), $P_2$ (quaternion conjugation), and
  $P_3$ (axis-flip) already established in the algebra, together with a
  corresponding action on the spacetime and complex-time arguments.
  A canonical choice of time-reversal convention $T_{\mathrm{UBT}}$ is stated
  and motivated.  The relation to $\psi$-parity $P_\psi$ and to the
  $SU(2)_L$ chirality derivation (\texttt{chirality/step1\_psi\_parity.tex})
  is clarified.
  \end{abstract}
\fi

\section{Discrete Symmetry Transformations: Overview}
\label{sec:sym:overview}

\subsection{Starting Point: Algebra and Field}

The biquaternion algebra is
\begin{equation}
  \mathcal{B} = \mathbb{C} \otimes \mathbb{H},
  \label{eq:sym:B}
\end{equation}
with ordered basis
$\{1, \mathbf{I}, \mathbf{J}, \mathbf{K}, i, i\mathbf{I}, i\mathbf{J},
i\mathbf{K}\}$,
where $i \in \mathbb{C}$ is the central imaginary unit and
$\mathbf{I}, \mathbf{J}, \mathbf{K} \in \mathbb{H}$ are quaternion units.
Three commuting $\mathbb{R}$-linear involutions $P_1, P_2, P_3$ on
$\mathcal{B}$ are defined in
\texttt{canonical/algebra/involutions\_Z2xZ2xZ2.tex}:
\begin{align}
  P_1: &\quad i \mapsto -i,\quad \mathbf{I},\mathbf{J},\mathbf{K}
        \mapsto \mathbf{I},\mathbf{J},\mathbf{K}
        \quad \text{(complex conjugation)}, \label{eq:sym:P1} \\
  P_2: &\quad \mathbf{I},\mathbf{J},\mathbf{K} \mapsto
        -\mathbf{I},-\mathbf{J},-\mathbf{K},\quad i \mapsto i
        \quad \text{(quaternion conjugation)}, \label{eq:sym:P2} \\
  P_3: &\quad \mathbf{J},\mathbf{K} \mapsto -\mathbf{J},-\mathbf{K},
        \quad \mathbf{I},i \mapsto \mathbf{I},i
        \quad \text{(axis-flip around $\mathbf{I}$)}. \label{eq:sym:P3}
\end{align}

The fundamental field is $\Theta(q,\tau) \in \mathcal{B}$, where:
\begin{itemize}
  \item $q = q^0 + q^1\mathbf{I} + q^2\mathbf{J} + q^3\mathbf{K} \in \mathcal{B}$
    is the biquaternion coordinate encoding spacetime position;
  \item $\tau = t + i\psi \in \mathbb{C}$ is the complex time parameter,
    $t \in \mathbb{R}$ (real time), $\psi \in \mathbb{R}$ (imaginary time phase).
\end{itemize}

We write the spatial vector part of $q$ as
$\mathbf{q} := q^1\mathbf{I} + q^2\mathbf{J} + q^3\mathbf{K}$, so that
$q = q^0 + \mathbf{q}$ and the quaternion conjugate is
$\bar{q} = q^0 - \mathbf{q}$.

%──────────────────────────────────────────────────────────────────────────────
\section{Charge Conjugation $C$}
\label{sec:sym:C}
%──────────────────────────────────────────────────────────────────────────────

\subsection{Definition}

\begin{definition}[Charge Conjugation $C$]
\label{def:sym:C}
Charge conjugation $C$ acts on the field by complex conjugation of all
biquaternion components:
\begin{equation}
  C[\Theta](q,\tau) := P_1\!\bigl(\Theta(q,\tau)\bigr)
                     = \overline{\Theta(q,\tau)}^{\,\mathbb{C}},
  \label{eq:sym:C_field}
\end{equation}
where $P_1$ is the complex-conjugation involution \eqref{eq:sym:P1}.
Explicitly, if $\Theta = \sum_a \Theta_a e_a$ with
$\Theta_a \in \mathbb{C}$ and
$e_a \in \{1, \mathbf{I}, \mathbf{J}, \mathbf{K}\}$, then
\begin{equation}
  C[\Theta] = \sum_a \Theta_a^* e_a.
  \label{eq:sym:C_explicit}
\end{equation}
The complex-time parameter $\tau = t + i\psi$ transforms as
\begin{equation}
  C: \tau \mapsto \tau^* = t - i\psi.
  \label{eq:sym:C_tau}
\end{equation}
For gauge fields $A_\mu^a$ (valued in the Lie algebra), $C$ acts as
\begin{equation}
  C: A_\mu^a T_a \mapsto -A_\mu^a T_a^*,
  \label{eq:sym:C_gauge}
\end{equation}
which reverses the sign of all U(1) charges and conjugates the
representation matrices $T_a$.
\end{definition}

\subsection{Identification with Algebra Involution}

At the level of the algebra value, $C = P_1$.  The transformation of the
complex-time argument $\tau \to \tau^*$ is a direct consequence of
$P_1$ acting on the imaginary unit $i \in \mathbb{C}$: since $\tau$ is
constructed from $i$, charge conjugation maps $\tau \to t - i\psi$.

\begin{remark}[Open question: internal gauge phase]
\label{rem:sym:C_gauge_phase}
Whether $C$ is identified purely with $P_1$ (complex conjugation of the
scalar complex factor) or must additionally include a non-trivial action on
internal U(1) phases carried by $\Theta$ depends on the detailed
representation.
For Abelian U(1), $C$ is uniquely $P_1$ composed with $A_\mu \to -A_\mu$.
For non-Abelian gauge fields, the action on $T_a$ must be specified by
the relevant group representation.  The definition above is canonical for
the U(1) and U(1)$_Y$ sectors.
\end{remark}

\subsection{Relation to $\psi$-Circle Orientation}

From \texttt{chirality/step1\_psi\_parity.tex}, the $\psi$-parity operator
$P_\psi: \psi \to -\psi$ sends winding number $n \to -n$, thereby exchanging
matter ($n > 0$) and antimatter ($n < 0$) sectors.
Since $C$ maps $\tau \to \tau^* = t - i\psi$, it acts as
$\psi \to -\psi$ on the imaginary time argument:
\begin{equation}
  \boxed{C \supseteq P_\psi
  \text{ (charge conjugation includes the $\psi$-parity of complex time)}.}
  \label{eq:sym:C_Ppsi}
\end{equation}
This confirms that the $\psi$-circle orientation identified in the chirality
derivation is a consequence of $C$-symmetry acting on complex time.

%──────────────────────────────────────────────────────────────────────────────
\section{Spatial Parity $P$}
\label{sec:sym:P}
%──────────────────────────────────────────────────────────────────────────────

\subsection{Definition}

\begin{definition}[Spatial Parity $P$]
\label{def:sym:P}
Spatial parity $P$ acts by flipping all spatial directions while
leaving time unchanged:
\begin{equation}
  P[\Theta](q,\tau)
  := P_2\!\bigl(\Theta(\bar{q},\tau)\bigr),
  \label{eq:sym:P_field}
\end{equation}
where $\bar{q} = q^0 - \mathbf{q}$ is the quaternion conjugate of the
coordinate (flipping the spatial vector part), and $P_2$ is the quaternion
conjugation involution \eqref{eq:sym:P2} acting on the algebra value
of $\Theta$.

In component notation, if
$q = q^0 + q^i \mathbf{e}_i$ (with $\mathbf{e}_i \in \{\mathbf{I},
\mathbf{J}, \mathbf{K}\}$) encodes $(t, x^1, x^2, x^3)$, then
$\bar{q}$ encodes $(t, -x^1, -x^2, -x^3)$, i.e.,
\begin{equation}
  P: (t, x^1, x^2, x^3) \mapsto (t, -x^1, -x^2, -x^3).
  \label{eq:sym:P_spacetime}
\end{equation}
The complex-time parameter is parity-even:
\begin{equation}
  P: \tau = t + i\psi \mapsto \tau.
  \label{eq:sym:P_tau}
\end{equation}
\end{definition}

\subsection{Action on Gauge Fields}

Under $P$, spatial components of gauge fields (vectors) pick up a sign:
\begin{equation}
  P: A_0(x) \mapsto A_0(\bar{x}),
  \qquad
  P: A_i(x) \mapsto -A_i(\bar{x}),
  \label{eq:sym:P_gauge}
\end{equation}
while field-strength components transform as:
\begin{equation}
  P: F_{0i}(x) \mapsto -F_{0i}(\bar{x}),
  \qquad
  P: F_{ij}(x) \mapsto F_{ij}(\bar{x}).
  \label{eq:sym:P_Fmunu}
\end{equation}
Here $\bar{x} = (t, -\mathbf{x})$.

\subsection{Left- and Right-Handed Decomposition}

The involution $P$ separates left-handed and right-handed biquaternionic
components.  As established in
\texttt{canonical/algebra/involutions\_Z2xZ2xZ2.tex}, the $P_2$-eigenspaces
are:
\begin{align}
  V_+ &= \mathrm{span}_\mathbb{C}\{1\}
        \quad (P_2\text{-eigenvalue }+1,\text{ scalar/even sector}), \\
  V_- &= \mathrm{span}_\mathbb{C}\{\mathbf{I},\mathbf{J},\mathbf{K}\}
        \quad (P_2\text{-eigenvalue }-1,\text{ vector/odd sector}).
\end{align}
The projections
$\Theta_\pm = \tfrac{1}{2}(1 \pm P_2)\Theta$
correspond to parity-even ($\Theta_+$, right-handed)
and parity-odd ($\Theta_-$, left-handed) components.
This is the origin of chirality asymmetry: the weak interaction couples
exclusively to $\Theta_-$ (see \texttt{chirality/step1\_psi\_parity.tex}).

%──────────────────────────────────────────────────────────────────────────────
\section{Time Reversal $T$}
\label{sec:sym:T}
%──────────────────────────────────────────────────────────────────────────────

\subsection{Three Candidate Definitions}

In UBT, complex time $\tau = t + i\psi$ introduces an ambiguity in the
definition of time reversal that does not appear in standard QFT.
Three candidate transformations of $\tau$ arise naturally:

\begin{definition}[Time Reversal Variants]
\label{def:sym:T_variants}
\begin{align}
  T_1: &\quad \tau \mapsto -t + i\psi
        \quad \text{(real part reversed, imaginary part preserved)},
        \label{eq:sym:T1} \\
  T_2: &\quad \tau \mapsto -\tau = -t - i\psi
        \quad \text{(full complex negation)},
        \label{eq:sym:T2} \\
  T_3: &\quad \tau \mapsto -\tau^* = -t + i\psi
        \quad \text{(complex conjugate then negate $=$ antiunitary on $\tau$)}.
        \label{eq:sym:T3}
\end{align}
Note that $T_3$ coincides with $T_1$ as a map on $\tau$.
They differ when acting on the field value: $T_3$ additionally complex-conjugates
$\Theta$ (antiunitary part), while $T_1$ does not.
\end{definition}

\begin{remark}
$T_3$ is the antiunitary time reversal in the sense of Wigner:
$T_3[\Theta](q,\tau) = P_1(\Theta(q, -\tau^*))$,
i.e., $t \to -t$ with $\psi$ unchanged, combined with complex conjugation of
the field value.  This is the standard physics convention for antiunitary $T$.
\end{remark}

\subsection{Canonical UBT Convention}

\begin{definition}[Canonical Time Reversal $T_{\mathrm{UBT}}$]
\label{def:sym:T_UBT}
The canonical UBT time reversal is chosen to be $T_2$:
\begin{equation}
  T_{\mathrm{UBT}}[\Theta](q,\tau)
  := \Theta(q, -\tau)
  = \Theta(q, -t - i\psi).
  \label{eq:sym:T_UBT}
\end{equation}
This is the \emph{full complex time negation}: both the real time $t$ and the
imaginary time $\psi$ reverse sign simultaneously.
\end{definition}

\paragraph{Motivation for $T_2$.}
\begin{enumerate}
  \item \textbf{Complex-scalar consistency.}
    The biquaternion framework treats $\tau \in \mathbb{C}$ as a single
    complex entity.  The most natural negation on $\mathbb{C}$ that reverses
    the time direction is $\tau \to -\tau$, preserving the complex-linear
    structure.
  \item \textbf{CPT structure.}
    Under $T_2$, the composition $C \circ P \circ T_2$ acts as:
    $(q,\tau) \to (\bar{q}, -\tau)$ with
    $\Theta \to P_1(P_2(\Theta(q, -\tau)))$.
    This reduces to the standard $CPT$ theorem in the real-time limit
    $\psi \to 0$ (see Section~\ref{sec:sym:CPT}).
  \item \textbf{Consistency with three-generation structure.}
    The mode expansion $\Theta = \sum_n \Theta_n e^{in\psi/R_\psi}$ under
    $T_2$ maps $e^{in\psi} \to e^{-in\psi}$, i.e., $n \to -n$.
    Combined with $C$ (which already sends $\psi \to -\psi$), the combined
    $CT_2$ acts as identity on the $\psi$-winding structure, consistent with
    the matter/antimatter interpretation.
\end{enumerate}

\subsection{Action on Gauge Fields under $T_{\mathrm{UBT}}$}

Under real-time reversal $t \to -t$ (part of $T_2$):
\begin{equation}
  T_{\mathrm{UBT}}: A_0(t,\mathbf{x}) \mapsto A_0(-t,\mathbf{x}),
  \quad
  A_i(t,\mathbf{x}) \mapsto -A_i(-t,\mathbf{x}),
  \label{eq:sym:T_gauge}
\end{equation}
(vector potential is time-odd for spatial components, time-even for the
temporal component, in order to preserve $F_{\mu\nu}$ invariance).
The field strength transforms as:
\begin{equation}
  T_{\mathrm{UBT}}: F_{0i} \mapsto F_{0i},
  \qquad
  F_{ij} \mapsto F_{ij}.
  \label{eq:sym:T_Fmunu}
\end{equation}

\begin{gapbox}
\textbf{Gap S1 --- Antiunitary Structure of $T_{\mathrm{UBT}}$:}
The canonical choice $T_2$ is defined here as a unitary transformation
(no complex conjugation of field values).  In standard QFT, the physically
correct time reversal is antiunitary.  Whether $T_{\mathrm{UBT}}$ should
be composed with $P_1$ (making it $T_3$ in the biquaternion sense) depends on
the inner product chosen for the Hilbert space of $\Theta$ modes.
Resolving this requires specifying the canonical sesquilinear form on the
space of $\Theta$ fields.  \textbf{Priority: MEDIUM.}
\end{gapbox}

%──────────────────────────────────────────────────────────────────────────────
\section{CPT Composition}
\label{sec:sym:CPT}
%──────────────────────────────────────────────────────────────────────────────

\subsection{Definition}

\begin{definition}[$CPT$ Operator]
\label{def:sym:CPT}
The $CPT$ operator is defined as the composition:
\begin{equation}
  CPT := C \circ P \circ T_{\mathrm{UBT}}.
  \label{eq:sym:CPT}
\end{equation}
Applied to the field:
\begin{align}
  CPT[\Theta](q,\tau)
  &= C\bigl[P\bigl[T_2[\Theta]\bigr]\bigr](q,\tau) \notag \\
  &= P_1\Bigl(P_2\bigl(\Theta(\bar{q}, -\tau)\bigr)\Bigr) \notag \\
  &= \bigl(P_1 \circ P_2\bigr)\bigl(\Theta(\bar{q}, -\tau)\bigr).
  \label{eq:sym:CPT_explicit}
\end{align}
\end{definition}

\begin{lemma}[$CPT$ action on coordinates and algebra]
\label{lem:sym:CPT_action}
Under $CPT$:
\begin{enumerate}
  \item The spacetime coordinate transforms as $q \to \bar{q}$
    (spatial inversion: $\mathbf{q} \to -\mathbf{q}$).
  \item The complex time transforms as $\tau \to -\tau$
    (full complex time reversal).
  \item The field value transforms by $P_1 \circ P_2$:
    \begin{equation}
      (P_1 \circ P_2)(b) = -b
      \quad \forall\, b \in \{\mathbf{I},\mathbf{J},\mathbf{K}\},
      \qquad
      (P_1 \circ P_2)(i) = -i,
      \qquad
      (P_1 \circ P_2)(1) = 1.
      \label{eq:sym:CPT_algebra}
    \end{equation}
\end{enumerate}
\end{lemma}

\begin{proof}
Direct computation from \eqref{eq:sym:P1}--\eqref{eq:sym:P2}:
$P_1$ sends $i \to -i$, $\mathbf{I},\mathbf{J},\mathbf{K} \to \mathbf{I},\mathbf{J},\mathbf{K}$;
$P_2$ sends $\mathbf{I},\mathbf{J},\mathbf{K} \to -\mathbf{I},-\mathbf{J},-\mathbf{K}$;
their composition gives $i \to -i$, $\mathbf{I},\mathbf{J},\mathbf{K} \to -\mathbf{I},-\mathbf{J},-\mathbf{K}$,
confirming \eqref{eq:sym:CPT_algebra}.
The coordinate and time transformations follow directly from
Definitions~\ref{def:sym:P} and \ref{def:sym:T_UBT}.
\end{proof}

\subsection{Reduction to Standard $CPT$ in the Real Limit}

\begin{proposition}[$CPT$ in the $\psi \to 0$ limit]
\label{prop:sym:CPT_real_limit}
In the real-time limit $\psi \to 0$ (i.e., $\tau \to t$),
the $CPT$ transformation \eqref{eq:sym:CPT_explicit} reduces to:
\begin{equation}
  CPT[\Theta]\bigl((t,\mathbf{x}), t\bigr)
  = (P_1 \circ P_2)\bigl(\Theta((-t,-\mathbf{x}), -t)\bigr),
  \label{eq:sym:CPT_real}
\end{equation}
which corresponds to the standard $CPT$ transformation of QFT:
inversion of all four spacetime coordinates $(t,\mathbf{x}) \to (-t,-\mathbf{x})$
combined with charge conjugation (complex conjugation of the field value).
\end{proposition}

\begin{resultbox}
\textbf{Summary of Step 1.}
The three discrete symmetries are defined as:
\begin{itemize}
  \item $C = P_1$ on field values + $\tau \to \tau^*$: charge conjugation.
  \item $P = P_2$ on field values + $q \to \bar{q}$: spatial parity.
  \item $T_{\mathrm{UBT}}$: field unchanged + $\tau \to -\tau$:
    full complex time reversal.
  \item $CPT = (P_1 \circ P_2)$ on field values + $(q,\tau) \to (\bar{q},-\tau)$:
    full discrete symmetry, reduces to standard $CPT$ when $\psi \to 0$.
\end{itemize}
All three transformations are expressed as compositions of the algebraic
involutions $P_1, P_2$ already present in $\mathcal{B}$, confirming that the
discrete symmetry structure is intrinsic to the biquaternion algebra.
\end{resultbox}

%──────────────────────────────────────────────────────────────────────────────
\section{Summary Table of Involution Identifications}
\label{sec:sym:table}
%──────────────────────────────────────────────────────────────────────────────

\begin{table}[h]
\centering
\caption{Correspondence between discrete symmetries and biquaternion involutions.}
\label{tab:sym:involutions}
\begin{tabular}{@{}lllll@{}}
  \toprule
  Symmetry & Algebra action & Coord.\ action & Time action & Involution \\
  \midrule
  $C$ (charge conj.) & $P_1$: $i \to -i$ & none & $\tau \to \tau^*$ & $P_1$ \\
  $P$ (parity) & $P_2$: $\mathbf{I},\mathbf{J},\mathbf{K} \to -$ & $q \to \bar{q}$ & $\tau \to \tau$ & $P_2$ \\
  $T$ (time rev.) & none & none & $\tau \to -\tau$ & (none) \\
  $CP$ & $P_1 \circ P_2$ & $q \to \bar{q}$ & $\tau \to \tau^*$ & $P_1 \circ P_2$ \\
  $CPT$ & $P_1 \circ P_2$ & $q \to \bar{q}$ & $\tau \to -\tau$ & $P_1 \circ P_2$ \\
  $P_\psi$ (parity of $\psi$) & none & none & $\psi \to -\psi$ & (sub-op.\ of $C$) \\
  \bottomrule
\end{tabular}
\end{table}

\begin{insightbox}
The involution $P_3$ (axis-flip around $\mathbf{I}$) does not appear as a
fundamental discrete symmetry $C$, $P$, or $T$ in isolation.
It acts as an inner automorphism of $\mathcal{B}$ and may represent a discrete
isospin rotation or a sub-symmetry within the $SU(2)_L$ sector.
Its role in symmetry breaking is discussed in
\texttt{symmetry/step3\_breaking\_catalogue.tex}.
\end{insightbox}

\ifdefined\INCLUDEMODE
\else
  \end{document}
\fi
